文言实词的活用，是指某个文言实词（主要是名词、形容词）经常用作甲类词，但在一定的语言环境中，按照一定的语言习惯灵活用作乙类词。例如，蒲松龄在《狼》中写两条狼紧紧跟踪一个屠户的情况后写道：“少时，一狼径去，其一犬坐于前。”（见于部编语文初中第一册）这个“犬”字本是名词，但从上下文来看，是不会又跑出一条“狗”来的，因此就根据这个“犬”字处在动词“坐”之前，把它看成修饰动词“坐”的状语，意思是“象狗似地”，这就叫名词用作状语，是词类活用中的一种。如果不了解这种词类活用的现象，而按照现代汉语的一般用法把句子理解成“其一犬，坐于前”，是要闹笑话的。

我们要正确理解文言文句子的意思，必须掌握住文言实词在句子中的活用词义。而文言实词那么多，我们不可能一一记住它们在句子中活用的词义。怎么办呢？只有掌握词类活用的规律才行。

具体说来，以下一些关于词类活用的规律值得认真领会。

1. 名词用作动词的规律：

我们判断一个名词是不是活用作动词，应从整个句子的意思来考虑，同时还要注意它在句中的地位，以及它前后有哪些词类的词跟它相配合，跟它构成什么样的语法关系。
就一般情况说，名词处在代词前面时活用作动词，因为代词不受名词修饰。例如：
\begin{enumerate}[label=\textcircled{\arabic*}]
	\item 宦官惧其毁己，皆共目（看）之。（《张衡传》）
	\item 驴不胜怒，蹄（踢）之。（《黔之驴》）
	\item 徐庶见先主，先主器（器重）之。（《隆中对》）
\end{enumerate}

如果名词处在副词特别是否定副词后面时，一般也活用作动词，因为名词不受副词特别是否定副词修饰。例如：

\begin{enumerate}[label=\textcircled{\arabic*}]
	\item 小信未孚，神弗福（赐福，保佑）也。（《曹刿论战》）
	\item 上胡不法（效法）先王之法？（《吕氏春秋·察今》）
	\item 如曰今日当一切不事（做）事，守前所为而已，则非某之所敢知。（《答司马谏议书》）
\end{enumerate}

如果名词处在可以肯定的宾语之前时，一般也活用作动词，因为宾语前面必须有谓语。例如：

\begin{enumerate}[label=\textcircled{\arabic*}]
	\item 吾师（学）道也。（《师说》）
	\item 范增数目（递眼色）项王。（《鸿门宴》）
	\item 秦王与赵王会饮，令赵王鼓（奏）瑟。（《廉颇蔺相如列传》）
\end{enumerate}

如果一个句子中有两个名词连着使用，而它们既不是并列关系，也不是偏正关系，那么其中一定有一个名词已经活用作动词了。例如：

\begin{enumerate}[label=\textcircled{\arabic*}]
	\item 乃丹书帛曰：“陈胜王（陈胜为王）。”（《陈涉世家》）
	\item 左右欲刃相如（用刀杀相如），相如张目叱之，左右皆靡。（《廉颇蔺相如列传》）
	\item 陛下不能将兵（带兵），而善将将（统率将领），此乃信（韩信的自称）之所以为陛下禽（通“擒”）也。（《史记·淮阴侯列传》）
	\item 豪贼劫持，反接，布囊其口（用布袋子堵住他的嘴）。（《童区寄传》）
\end{enumerate}

如果能原动词如“能”、“可”、“肯”等后面只有名词，而没有另外的动词，这个名词就活用作动词。例如：

\begin{enumerate}[label=\textcircled{\arabic*}]
	\item 假舟楫者，非能水（游泳）也，而绝江河。（《劝学》）
	\item 秦地可尽王（为王）也。（《鸿门宴》）
	\item 人有百口，口有百舌，不能名（说出）其一处也。（《口技》）
	\item 三岁贯女（通“汝”），莫我肯德（施恩）。（《硕鼠》）
\end{enumerate}

当然，一个名词处于句子谓语的地位，肯定活用作动词。例如：

\begin{enumerate}[label=\textcircled{\arabic*}]
	\item 今王与百姓同乐，则王（称王，统一天下）矣。（《庄暴见孟子》）
	\item 唐浮图慧褒始舍（居住）于其址。（《游褒禅山记》）
	\item 江水又东（向东流）。（《水经注·江水》）
\end{enumerate}

另外，如果“而”字连接的两个部分，其中一部分是动词或动词性词组，而另一部分是名词的话，这个名词一般也活用作动词。例如：

\begin{enumerate}[label=\textcircled{\arabic*}]
	\item 汉败楚，楚以故不能过荥阳而西（向西进）。（《史记·项羽本纪》）
	\item 先生且喜且愕，舍狼而前（向前去）。（《中山狼传》）
\end{enumerate}

2.形容词用作动词的规律：

如果一个文言句子，把形容词直接放在宾语前面，作为句子的谓语，这个形容词就活用作动词。例如：

\begin{enumerate}[label=\textcircled{\arabic*}]
	\item 西人长（擅长）火器而短（不擅长）技击。（《冯婉贞》）
	\item 尔安敢轻（轻视）吾射！（《卖油翁》）
	\item 京中有善（擅长）口技者。（《口技》）
	\item 大王必欲急（逼迫）臣，臣头今与璧俱碎于柱矣！”（《廉颇蔺相如列传》）
	\item 卒使上官大夫短（诋毁）屈原于顷襄王。（《屈原列传》）
\end{enumerate}

3. 动词、形容词用作名词的规律：

如果两个动词直接连用，往往后面的动词用作名词，作为前面动词的宾语。例如：

追亡（逃亡者）逐北（败北者），伏尸百万。（《过秦论》）

如果形容词放在动词之后作动词的宾语，这个形容词则活用作名词。例如：

\begin{enumerate}[label=\textcircled{\arabic*}]
	\item 当横行天下，为汉家除残（残暴者）去秽（污秽者）。 （《赤壁之战》）
	\item 将军披坚（坚固的盔甲）执锐（锐利的兵器），伐无道，诛暴秦······（《陈涉世家》）
\end{enumerate}

4. 名词、动词用作状语的情况：

文言文中的普通名词，有时可以直接修饰动词，用作状语。它同主谓结构不同之点是：名词用做主语时，后面的动词是陈述这个主语怎么样，做什么；而名词用作状语时，则是修饰它后面的动词象什么似的或者在怎样的方式与情况下发出这个动作。这个区别，只要把句子的上下文意联系起来看，是很容易辨明的。例如：

\begin{enumerate}[label=\textcircled{\arabic*}]
	\item 将不胜其忿而蚁（象蚂蚁一样）附之。（《谋攻》）
	\item （狼）猬缩蠖屈，蛇盘龟息，以听命先生。（《中山狼传》）
	
	（加点的名词均用作后面动词的状语，有“象······一样”的意思。）
	
	\item 项伯亦拔剑起舞，常以身翼（如鸟张翼似地）蔽沛公，庄不得击。（《鸿门宴》）
	齐将田忌善而客待之。（《史记·孙子吴起列传》）（齐国的将军田忌很赞许他，并且把他当作上宾来接待。“善”是形容词用作动词，“客”是名词用作状语，修饰动词“待”，意思是“象对待客人一样”。）
	\item 夫以秦王之威，而相如廷（在朝廷上）叱之。（《廉颇蔺相如列传》）
	\item 事不目（用眼睛）见耳（用耳朵）闻，而臆断其有无，可乎？（《石钟山记》）
	\item 良庖岁（每年）更刀，割也；族庖月（每月）更刀，折也。（《庖丁解牛》）
	\item 而乡邻之生日（一天天地）蹙。（《捕蛇者说》）
\end{enumerate}

由上可见，普通名词用作状语，在表示比喻时，有“象······一样”的意思；在表示对待人的态度时，有“象对待······一样”的意思；在表示处所和工具时，可以加上“在”、“用”等适当的介词去理解；如果时间名词用作状语，那么，“岁”、“月”、“日”放在具有行动性的动词前面就有“每年”、“每月”、“每天”的意思，表示行动的频数或经常；“日”若放在动词或形容词前面，当“一天天地”讲，表示情况的逐渐发展。

动词有时也可以用作状语。这个活用作状语的动词，一般只限于不及物动词，后面往往紧接着被修饰的动词。例如：

\begin{enumerate}[label=\textcircled{\arabic*}]
	\item 儿惧，啼（哭着）告母。（《促织》）
	\item 邻人京城氏之孀妻有遗男，始龀，跳（蹦蹦跳跳地）往助之。（《愚公移山》）
	\item 往来（走过来走过去地）视之，觉无异能者。（《黔之驴》）
\end{enumerate}

5. 动词、形容词、名词用作使动词的规律：
所谓使动用法，就是使宾语所代表的人或事物施行某个动作，或具有某种性质状态，或成为某种事物。即是说，在文言文中，动词、形容词或名词如果用作使动词的话，它们对于宾语都含有“使它怎样”的意思。
例如动词的使动用法：

\begin{enumerate}[label=\textcircled{\arabic*}]
	\item 坐定，公子从车骑（使车骑跟着），虚左，自迎夷门侯生。（《信陵君窃符救赵》）
	\item 外连横而斗诸侯（使诸侯相斗）。（《过秦论》）
	\item 令女居其上，浮之（使之浮）河中。始浮，行数十里乃没。（《西门豹治邺》）
	\item 广故数言欲亡，忿恚尉（使尉忿恚）。（《陈涉世家》）
	\item 项伯杀人，臣活之（使之活）。（《鸿门宴》）
\end{enumerate}

这就说明：只要动宾之间不是支配和被支配的关系，宾语不是动作的承受者，而是主语使宾语所代表的事物施行这个动作，这个动词就用作使动词。这种动词大多是通常不带宾语的不及物动词（如“忿恚”、“怒”、“活”、“死”。“入”、“来”等等）。

再如形容词的使动用法：

\begin{enumerate}[label=\textcircled{\arabic*}]
	\item 乃出图书，空囊橐（使口袋空着），徐徐焉实狼其中。（《中山狼传》）
	\item 攻邯郸，以广地尊名（使地广，使名尊）。（《战国策·南辕北辙》）
	\item 臣请完璧（使璧完美无缺）归赵。（《廉颇蔺相如列传》）
	\item 春风又绿（使江南岸绿）江南岸，明月何时照我还。（《泊船瓜洲》）
	\item 于是废先王之道，焚百家之言，以愚黔首（使百姓愚）。（《过秦论》）
\end{enumerate}

这就说明：如果形容词处在带着宾语的及物动词的位置上，表示主语“使”宾语所代表的事物具有这个形容词所表示的性质，这个形容词就用作使动词。

又如名词的使动用法：

\begin{enumerate}[label=\textcircled{\arabic*}]
	\item 然而而腊之（使之成为干肉）以为饵，可以已大风······（《捕蛇者说》）
	\item 纵江东父老怜而王我（使我为王），我何面目见之。（《史记·项羽本纪》）
	\item 父曰：“履我（给我穿鞋）！”良业为取履，因长跪履之（给他穿鞋）。（《史记·留侯世家》）
\end{enumerate}

这就说明：如果名词处在带着宾语的及物动词的位置上，表示主语“使”宾语所代表的事物成为这个名词所代表的事物，或具有这个名词的属性，这个名词就用作使动词。

判别是不是使动用法，要根据本句的具体内容和上下文。如“饮马长江”，根据“马”和“饮”的关系，就可以断定只能是“使马饮”。又如《中山狼传》中的“是狼为虞人所窘，求救于我，我实生之”这句话，根据句意和上下文，“生之”只能是“使它活”的意思。

6.形容词、名词用作意动词的规律：

所谓意动用法，就是主观上认为宾语所代表的人或事物具有某种性质状态，或者成为某种事物。在文言文中，形容词和名词都有意动用法，它们对于其宾语都含有“认为它怎样”或者“将它看作什么”的意思。

例如形容词用作意动词：

\begin{enumerate}[label=\textcircled{\arabic*}]
	\item 大将军邓骘奇其才（以其才为奇）。（《张衡传》）
	\item 吾妻之美我（以我为美）者，私我也。（《邹忌讽齐王纳谏》）
	\item 成以其小，劣之（以之为劣）。（《促织》）
	\item 得鱼腹中书，固以怪之（以之为怪）矣。（《陈涉世家》）
	\item 孔子登东山而小鲁（认为鲁小），登泰山而小天下（认为天下小）。（《孟子·尽心章句上》）
	\item 再如名词用作意动词：
	\item 吾从而师之（以之为师）。（《师说》）
	\item 惟大辟无可要（要挟），然犹质其首（以其首为质。“质”是“抵押品”的意思，名词用作意动词）。（《狱中杂记》）
	\item 邑人奇之，稍稍宾客其父（把其父看成宾客），或以钱币乞之。（《伤仲永》）
	\item 孟尝君客我（把我当作客人）。（《战国策·冯谖客孟尝君》）
	\item 粪土当年万户侯（把当年万户侯看作粪土）。（毛泽东《沁园春·长沙》）
\end{enumerate}

一般说来，代词前面的形容词用如使动词或意动词，同代词前面的名词用如使动词或意动词的道理相同，因为代词既不被名词所修饰，也不被形容词所修饰，代词前面的名词或形容词只能用如动词（使动、意动）。至于同一个形容词往往既可用如使动，又可用如意动，怎么区分呢？这只有靠分析上下文。例如同样一个“苦”字，是形容词，在“天将降大任于斯人也，必先苦其心志，劳其筋骨，饿其体肤，空乏其身”（《孟子·告子下》）中用如使动词，而在“天下苦秦久矣”（《陈涉世家》）中则用如意动词。

7. 数词用作动词的规律：

在古汉语中，数词有时充当句子的谓语，此时这些数词不再表示特定的数目，而用作动词了。经常被用作动词的数词是“一”和“二”，有时“二”和“三”连用，也能用作动词；有时别的数词也能用作动词表示人或事物数量上的变化，或者作句子的谓语，对主语起说明作用。例如：

\begin{enumerate}[label=\textcircled{\arabic*}]
	\item 六王毕，四海一（统一）。（《阿房宫赋》）
	\item 士也罔极（无常），二三（不专一）其德。（《诗经·卫风·氓》）
	\item 此三子者······与臣而将四（成为四个）矣。（《唐睢不辱使命》）
	\item 范增······举所佩玉块以示之者三（有三次）。（《鸿门宴》）
\end{enumerate}
\newpage